\section*{Independent reconstruction of the retained high-degree GU reserve}

All logarithms are base two.  This review audits the exact claim
\texttt{B20-A097-HIGH-DEGREE-RESERVE-v1}, whose frozen candidate is
\[
 G(f)\le4I(f)+{3\over4}U(f)                           \tag{1}
\]
for every genuine Boolean function on every finite uniform cube.

\paragraph{Audited source line.}
The complete original row-cut proof, before its final relaxation, gives
\[
 G\le D_{\ge2}+{1\over2}V_1+{2\over\ln2}(M_1+W_2),    \tag{2}
\]
where
\[
 M_1=\sum_i\min\{p(\varnothing),p(\{i\})\},\qquad
 W_d=\sum_{|S|=d}p(S),\qquad I=\sum_{d\ge1}dW_d.
\]
Equation (2), rather than its weaker version with $I$ replacing
$M_1+W_2$, is the necessary original for this audit.

\paragraph{Reserve identity and sign.}
Define
\[
 Q_{\ge3}=I-M_1-W_2.
\]
Because $M_1\le W_1$,
\[
 \begin{aligned}
 Q_{\ge3}
 &=W_1+2W_2+\sum_{d\ge3}dW_d-M_1-W_2\\
 &=(W_1-M_1)+W_2+\sum_{d\ge3}dW_d\\
 &\ge W_2+3W_{\ge3}.                                  \tag{3}
 \end{aligned}
\]
Thus $Q_{\ge3}\ge0$ always and is strictly positive whenever the
Fourier degree is at least three.  Substitution in (2) has the sign
\[
 \boxed{
 G\le D_{\ge2}+{1\over2}V_1
       +{2\over\ln2}(I-Q_{\ge3}).}                    \tag{4}
\]
No KL row, signed correlation, normalization term, or influence payment
is altered by this algebraic substitution.

\paragraph{Exact implication of the new leaf.}
Suppose, only on tensor-prime positive-residual Boolean functions of
degree at least three, that
\[
 D_{\ge2}+{1\over2}V_1
 \le\left(4-{2\over\ln2}\right)I+{3\over4}U
       +{2\over\ln2}Q_{\ge3}.                         \tag{5}
\]
Adding the remaining term from (4) cancels $Q_{\ge3}$ exactly and gives
(1).  The source-pinned tensor additivity localizes a positive residual
to a tensor-prime factor.  The independently proved degree-two theorem
$G\le2I$ makes the frozen residual strictly negative for every
nonconstant degree-at-most-two function.  Therefore (5) is sufficient
for universal (1).  Replacing $Q_{\ge3}$ in (5) by the smaller
$W_2+3W_{\ge3}$ gives a stronger, hence still sufficient, hypothesis.

This is a strict quantitative reduction on the remaining domain because
$Q_{\ge3}>0$ there.  It is not a proof of (5): high-degree mass can be
arbitrarily small, and no audited source bounds the remaining signed
higher-order KL excess by this reserve.

\paragraph{Independent GU-to-FEI recurrence.}
For completeness, assume a universal pair
\[
 G(f)\le aI(f)+bU(f),\qquad 0\le b<1.                 \tag{6}
\]
Set $C=a/(1-b)$.  For an $n$-variable function, exact restriction
identities give
\[
 \sum_{i=1}^n {U(f_{i,+})+U(f_{i,-})\over2}=nU(f)-G(f),
\]
\[
 \sum_{i=1}^n {I(f_{i,+})+I(f_{i,-})\over2}=(n-1)I(f). \tag{7}
\]
Using the induction hypothesis $U(f_{i,\pm})\le CI(f_{i,\pm})$ and
(6),
\[
 (n-b)U-aI\le nU-G\le C(n-1)I.
\]
Since
\[
 a+C(n-1)=C(n-b),
\]
one obtains $U\le CI$.  Constants give the zero base case separately.
The source reduction
\[
 H\le2(I-\operatorname{Var}f)+2U\le2I+2U
\]
then yields
\[
 H\le\left(2+{2a\over1-b}\right)I.                  \tag{8}
\]
For $(a,b)=(4,3/4)$ this is $U\le16I$ and $H\le34I$.

\paragraph{Verdict and first unsupported line.}
The claim \texttt{B20-A097-HIGH-DEGREE-RESERVE-v1} passes at its exact
reduction scope.  No material repair was made, so this review remains
independent of the proof's authorship.  Universal GU and FEI remain open.
The first unsupported line is exactly (5), with all its quantifiers:
it must hold for every tensor-prime, positive-residual genuine Boolean
function of arbitrary finite dimension and Fourier degree at least three.
